\documentclass[12pt]{article}
\usepackage[utf8]{inputenc}
\usepackage[T1]{fontenc}
\usepackage{hyperref}
\title{Kuwema Language Resources for K-12 Curriculum Development}
\author{Compiled by Claude Sonnet 5.5 (claude-sonnet-5-5) for Andr\'as Kornai}
\date{}
\begin{document}
\maketitle

\section{Introduction}
Kuwema (also Tyaraity, Tyeraity, Djerait) has the ISO 639-3 code woa and the Glottocode tyar1236. Glottolog lists it as dormant, with about 9 speakers recorded. The name is ambiguous. Kuwema is treated as a variety closely tied to Malak-Malak (Northern Daly), and Kandjerramalh, a separate language, is also sometimes called Kuwema. Search found very little material specific to Kuwema. This report does not invent anything beyond what appeared in search results.

\section{Primary Educational Resources}
No primers, textbooks or school programmes for Kuwema were found. The language is not known to be taught in schools.

\section{Linguistic and Reference Materials}
Glottolog has a languoid page for Kuwema: \url{https://glottolog.org/resource/languoid/id/tyar1236}.

The Endangered Languages Project has a bibliography page that may relate to this language: \url{https://endangeredlanguages.com/lang/3329/bibliography}. I did not confirm which language record the page covers.

Wikipedia's Malak-Malak article discusses the Kuwema/Tyaraity variety and its relationship to Malak-Malak: \url{https://en.wikipedia.org/wiki/Malak-Malak_language}.

Wikipedia's Kandjerramalh article is relevant for the name overlap, to avoid confusion: \url{https://en.wikipedia.org/wiki/Kandjerramalh_language}.

\section{Religious and Community Resources}
None found.

\section{Digital and Audio Resources}
Wikitongues has a page for the code woa: \url{https://wikitongues.org/languages/woa/}. I did not verify whether it holds any recordings.

\section{Curriculum Development Context}
Kuwema is dormant, with few or no fluent speakers. A heritage curriculum would have to rely on archival material and on work with the closely related Malak-Malak. Any such work should be done with the relevant Aboriginal community. No community organisation was identified in the searches.
\end{document}